\section{Metodologia Expositiva}
Lima (1973, apud \cite{andreataAulaExpositivaPaulo2019}, p. 444) contextualiza as origens do método expositivo, processo primordial de comunicação oral; como anterior à descoberta da escrita e da imprensa. A exemplo do quão antiga é a ferramenta, propõe uma ambientação no contexto do ensino universitário medieval, que dispunha apenas deste recurso para ensinar os alunos. Em particular, o cenário para o docente era o seguinte:

\begin{quote}
    A inexistência de livros levava-o a repetir (recitações-leitura) o texto tantas vezes quantas necessárias para que os alunos o decorassem (hipótese em que se poderia supor que os alunos podiam ser até analfabetos)(Lima, 1973, p. 444).
\end{quote}

Primeiro pedagogo da história, o filósofo grego Platão (427-347 a.C.) fazia uso do ensino oral em suas discussões, que - para Sócrates (469/470-399 a.C.), mestre de Platão; e os sofistas - tinha papel não só educativo, mas de produção da própria educação; simultâneamente visando a formação do indivíduo, possível de reconhecer pela aquisição de uma "competência discursiva". Gauthier; Tardif (2013, apud \cite{andreataAulaExpositivaPaulo2019}, p. 36) ainda complementa dizendo que, para os sofistas, ser instruído é "saber falar" e "saber argumentar" em público segundo as regras pragmáticas da retórica.

Andreata (2019) dialoga com Mattos (1959, apud \cite{andreataAulaExpositivaPaulo2019}),
Nérici (1968, apud \cite{andreataAulaExpositivaPaulo2019}) e Piletti (1987, apud \cite{andreataAulaExpositivaPaulo2019}) de forma a definir a metodologia expositiva. Em resumo, pode-se entendê-la como técnica de exposição oral, por parte do docente, de um determinado assunto, valendo-se de uma apresentação logicamente estruturada. Nela, o aluno tem papel de ouvinte, sendo esse responsável por assimilar o conteúdo ministrado pelo docente.

Há, no entanto - desde a Antiga Grécia - um grande risco para essa técnica, que mais tarde será reatado por Roma e a Idade Média: o uso majoritário de discursos preparados e repetição de memória contribuiam para um saber superficial e pretencioso, como destaca Gal (1989, apud \cite{andreataAulaExpositivaPaulo2019}, p. 32).

As principais críticas à aula expositiva são centradas na passividade do aluno, que é relegado no seu próprio processo de aprendizagem. Baldan (2011) \cite{baldanRepresentacaoAtoEnsinar2011} descreve o ato de ensinar, no sentido tradicional, da seguinte forma:
\begin{quote}
    [...] processo mecânico, no qual o professor detém o conhecimento e o aluno assume a função passiva de assimilação deste conteúdo. Os métodos são mecânicos, com repetição e cópia, aula expositiva e arguição oral, com o uso de instrumentos disciplinadores e reguladores do comportamento (Baldan, 2011).
\end{quote}

Outros autores, como de Lima (2022) \cite{delimaCRITICAMETODOLOGIATRADICIONAL} acrescentam dizendo que o problema da metodologia tradiconal é a falta de interação entre o sujeito e o objeto, a falta de diálogo entre o professor e o aluno; pois muitas vezes o que é exposto não faz dimensão alguma com a realidade do aluno presente. Isso ocorre porque há uma estrutura inflexível e unilateral, em que o docente "despeja" o saber e pouco faz para adaptar essa dinâmica às peculiaridades do aluno. Tal realidade se faz mais desafiadora no contexto do Ensino à Diastância (EAD), no qual há maior distânciamento entre sujeito e objeto; compromentendo a dinâmica já comprovada ineficaz em ambientes convencionais (sala de aula)(Vasconcellos, 1992 apud \cite{delimaCRITICAMETODOLOGIATRADICIONAL}).

Isto posto, é uma técnica já alardeada por inúmeros autores no que diz respeito à sua eficácia, inclusive na contemporâniedade, na qual há pleno aparato tecnológico para identificação de gargalos e solução dos problemas crônicos da aula expositiva. Adicionalmente, deve-se levar em conta sua popularidade atemporal; que convém o título de metodologia mais  usada em território brasileiro (Nérici (1968), apud \cite{andreataAulaExpositivaPaulo2019}.p.258).

\section{Ensino Remoto e EAD}
"Sistema tecnológico de comunicação bidirecional, que pode ser massivo, baseado em uma ação sistemática e conjunta de recursos didáticos e o apoio de tutoria, que, separados fisicamente dos estudantes, propiciam a esses uma aprendizagem independente" Arteio(2001, apud \cite{rochaMetodologiasAprendizagensNo2022}, p.39), o Ensino à Diastância (EAD) é uma modalidade de ensino crescente, em especial para o nível superior; com 44\% das matrículas brasileiras (2020) \cite{rochaMetodologiasAprendizagensNo2022}.  Dados do Instituto Nacional de Estudos e Pesquisas Educacionais Anísio Teixeira (INEP) apontam para um crescimento de 378,9\% no número de matrículas em graduações à distância, no período de 2009 à 2019. Dessa forma, é inegável que o EAD faz-se cada vez mais presente nos cursos superiores brasileiros. 

No entanto, mesmo sujeita às inovações proporcionadas pelas tecnologias da informação e comunicação (TIC), Suhre Ribeiro(2010, apud \cite{schneiderSalaAulaInvertida2013}, p.31) apontam para uma transposição da sala de aula tradicional para o contexto remoto, apenas. E isso é extremamente prejudicial a um modelo que permite atender às diferentes necessidades de aprendizado dos alunos, como destaca Belloni (2002)\cite{belloniEnsaioSobreEducacao2002}: "A EAD permite que se explore a diversidade inerente aos grupos sociais, produzindo efetividade na experiência educativa". Trazer a metodologia expositiva, cuja comunicação é unilateral e centrada no professor é conceitualmente um erro, no que diz respeito ao ensino remoto. 
\begin{quote}
    A EAD tem como principal característica, reduzir distâncias entre docentes e discentes (bem como demais interagentes), sejam elas espaciais, culturais, sociais e temporais (\cite{belloniEnsaioSobreEducacao2002}).
\end{quote}



\section{Atenção e Engajamento no Processo de Aprendizagem}

\section{Tecnologia no Monitoramento da Aprendizagem}

\section{Engajamento e Supervisão do Processo de Ensino Online}